\documentclass{beamer}

\usepackage{amsmath}
\usepackage{amsfonts}
\usepackage{amssymb}
\usepackage{bm}
\usepackage{graphicx}

\usetheme{Madrid}
\usecolortheme{default}

\title{2D Fourier Phase-Shift Displacement Est.}
\subtitle{GPR Time-Lapse Monitoring}
\author{Conner van Kooten}
\institute{Thesis Meeting}
\date{9 June 2026}

% ─────────────────────────────────────────────────────────────────────────────
\begin{document}

\maketitle

% ─── Motivation frame ────────────────────────────────────────────────────────
\begin{frame}{Motivation: Linearity of Gazdag Migration}

\begin{columns}[c]
  \begin{column}{0.48\textwidth}
    \centering
    \includegraphics[width=\linewidth]{Gazdag_migration_image.png}\\[2pt]
    {\small Migrate baseline \& monitor,\\then subtract}
  \end{column}
  \hfill
  \begin{column}{0.48\textwidth}
    \centering
    \includegraphics[width=\linewidth]{Gazdag_migration_timelapse_image.png}\\[2pt]
    {\small Subtract B-scans first,\\then migrate the difference}
  \end{column}
\end{columns}

\bigskip
\begin{block}{Key observation}
  Because Gazdag migration is a \textbf{linear} operation, both workflows
  produce identical results:
  \[
    \mathcal{M}\{b\} - \mathcal{M}\{m\} \;=\; \mathcal{M}\{b - m\}
  \]
\end{block}

\end{frame}

% ─── Migration gallery ───────────────────────────────────────────────────────
\begin{frame}{Migration Methods — Gallery}

\begin{center}
  \includegraphics[width=\linewidth,height=0.80\textheight,keepaspectratio]{Migration_galary.png}
\end{center}

\end{frame}

% ─── Lateral shift experiment ────────────────────────────────────────────────
\begin{frame}{Lateral Shift Experiment}

\begin{columns}[c]
  \begin{column}{0.32\textwidth}
    \centering
    \includegraphics[width=\linewidth]{Lateral_shift_test_galary_0p25_lambda.png}\\[2pt]
    {\small $\Delta x = \tfrac{1}{4}\lambda$}
  \end{column}
  \begin{column}{0.32\textwidth}
    \centering
    \includegraphics[width=\linewidth]{Lateral_shift_test_galary_0p125_lambda.png}\\[2pt]
    {\small $\Delta x = \tfrac{1}{8}\lambda$}
  \end{column}
  \begin{column}{0.32\textwidth}
    \centering
    \includegraphics[width=\linewidth]{Lateral_shift_test_galary_0p0625_lambda.png}\\[2pt]
    {\small $\Delta x = \tfrac{1}{16}\lambda$}
  \end{column}
\end{columns}

\end{frame}

% ─── PSF phase change ────────────────────────────────────────────────────────
\begin{frame}{Point Spread Functions \& Phase Change}

\begin{columns}[c]
  \begin{column}{0.32\textwidth}
    \centering
    \includegraphics[width=\linewidth]{Lateral_shift_test_psf_0p25_lambda.png}\\[2pt]
    {\small $\Delta x = \tfrac{1}{4}\lambda$}
  \end{column}
  \begin{column}{0.32\textwidth}
    \centering
    \includegraphics[width=\linewidth]{Lateral_shift_test_psf_0p125_lambda.png}\\[2pt]
    {\small $\Delta x = \tfrac{1}{8}\lambda$}
  \end{column}
  \begin{column}{0.32\textwidth}
    \centering
    \includegraphics[width=\linewidth]{Lateral_shift_test_psf_0p0625_lambda.png}\\[2pt]
    {\small $\Delta x = \tfrac{1}{16}\lambda$}
  \end{column}
\end{columns}

\end{frame}

% ─── PSF phase change (zoomed) ───────────────────────────────────────────────
\begin{frame}{Point Spread Functions \& Phase Change — Zoomed}

\begin{columns}[c]
  \begin{column}{0.32\textwidth}
    \centering
    \includegraphics[width=\linewidth,height=0.78\textheight,keepaspectratio]{Lateral_shift_test_psf_0p25_lambda_zoomed.png}\\[2pt]
    {\small $\Delta x = \tfrac{1}{4}\lambda$}
  \end{column}
  \begin{column}{0.32\textwidth}
    \centering
    \includegraphics[width=\linewidth,height=0.78\textheight,keepaspectratio]{Lateral_shift_test_psf_0p125_lambda_zoomed.png}\\[2pt]
    {\small $\Delta x = \tfrac{1}{8}\lambda$}
  \end{column}
  \begin{column}{0.32\textwidth}
    \centering
    \includegraphics[width=\linewidth,height=0.78\textheight,keepaspectratio]{Lateral_shift_test_psf_0p0625_lambda_zoomed.png}\\[2pt]
    {\small $\Delta x = \tfrac{1}{16}\lambda$}
  \end{column}
\end{columns}

\end{frame}

% ─── dΦ/dx gradient ──────────────────────────────────────────────────────────
\begin{frame}{Phase Gradient $\partial\Phi/\partial k_x$ at Migration Apex}

\begin{columns}[c]
  \begin{column}{0.32\textwidth}
    \centering
    \includegraphics[width=\linewidth]{dphi_dx_0p25_lambda.png}\\[2pt]
    {\small $\Delta x = \tfrac{1}{4}\lambda$}
  \end{column}
  \begin{column}{0.32\textwidth}
    \centering
    \includegraphics[width=\linewidth]{dphi_dx_0p125_lambda.png}\\[2pt]
    {\small $\Delta x = \tfrac{1}{8}\lambda$}
  \end{column}
  \begin{column}{0.32\textwidth}
    \centering
    \includegraphics[width=\linewidth]{dphi_dx_0p0625_lambda.png}\\[2pt]
    {\small $\Delta x = \tfrac{1}{16}\lambda$}
  \end{column}
\end{columns}

\end{frame}

% ─── Vertical shift experiment ───────────────────────────────────────────────
\begin{frame}{Vertical Shift Experiment}

\begin{columns}[c]
  \begin{column}{0.47\textwidth}
    \centering
    \includegraphics[width=\linewidth,height=0.36\textheight,keepaspectratio]{Vertical_shift_test_galary_0p25_lambda.png}\\[1pt]
    {\small $\Delta z = \tfrac{1}{4}\lambda$}\\[4pt]
    \includegraphics[width=\linewidth,height=0.36\textheight,keepaspectratio]{Vertical_shift_test_galary_0p0625_lambda.png}\\[1pt]
    {\small $\Delta z = \tfrac{1}{16}\lambda$}
  \end{column}
  \begin{column}{0.47\textwidth}
    \centering
    \includegraphics[width=\linewidth,height=0.36\textheight,keepaspectratio]{Vertical_shift_test_galary_0p125_lambda.png}\\[1pt]
    {\small $\Delta z = \tfrac{1}{8}\lambda$}\\[4pt]
    \includegraphics[width=\linewidth,height=0.36\textheight,keepaspectratio]{Vertical_shift_test_galary_0p03125_lambda.png}\\[1pt]
    {\small $\Delta z = \tfrac{1}{32}\lambda$}
  \end{column}
\end{columns}

\end{frame}

% ─── Vertical PSF phase change (zoomed) ──────────────────────────────────────
\begin{frame}{Vertical Shift — PSFs \& Phase Change}

\begin{columns}[c]
  \begin{column}{0.47\textwidth}
    \centering
    \includegraphics[width=\linewidth,height=0.36\textheight,keepaspectratio]{Vertical_shift_test_psf_0p25_lambda_zoomed.png}\\[1pt]
    {\small $\Delta z = \tfrac{1}{4}\lambda$}\\[4pt]
    \includegraphics[width=\linewidth,height=0.36\textheight,keepaspectratio]{Vertical_shift_test_psf_0p0625_lambda_zoomed.png}\\[1pt]
    {\small $\Delta z = \tfrac{1}{16}\lambda$}
  \end{column}
  \begin{column}{0.47\textwidth}
    \centering
    \includegraphics[width=\linewidth,height=0.36\textheight,keepaspectratio]{Vertical_shift_test_psf_0p125_lambda_zoomed.png}\\[1pt]
    {\small $\Delta z = \tfrac{1}{8}\lambda$}\\[4pt]
    \includegraphics[width=\linewidth,height=0.36\textheight,keepaspectratio]{Vertical_shift_test_psf_0p03125_lambda_zoomed.png}\\[1pt]
    {\small $\Delta z = \tfrac{1}{32}\lambda$}
  \end{column}
\end{columns}

\end{frame}

% ─── Frame 1 ─────────────────────────────────────────────────────────────────
\begin{frame}{1. The 2D Fourier Shift Theorem}

Let the baseline migrated GPR image be $b(z, x)$.
When a point scatterer moves, the monitor image is a translated copy:
\[
  m(z, x) = b(z - \Delta z,\; x - \Delta x)
\]

\bigskip
The 2D Fourier transform of the baseline image is:
\[
  B(k_z, k_x)
  = \int_{-\infty}^{\infty}\!\int_{-\infty}^{\infty}
    b(z, x)\,e^{-j(k_z z + k_x x)}\,dz\,dx
\]

\bigskip
By the \textbf{Fourier Shift Theorem}, shifting in space becomes a linear phase
rotation in the frequency domain:
\[
  M(k_z, k_x) = B(k_z, k_x) \cdot e^{-j(k_z \Delta z + k_x \Delta x)}
\]

\medskip
\begin{block}{Key implication}
  The amplitude spectrum is unchanged ($|M| = |B|$), but every frequency
  coordinate $(k_z, k_x)$ receives a unique phase shift.
\end{block}

\end{frame}

% ─── Frame 2 ─────────────────────────────────────────────────────────────────
\begin{frame}{2. Isolating the Phase Plane — The Cross-Spectrum}

The complex cross-spectrum $XS$ is formed by multiplying the baseline spectrum
by the complex conjugate of the monitor spectrum:
\[
  XS(k_z, k_x) = B(k_z, k_x) \cdot M^{*}(k_z, k_x)
\]

Substituting the shift equation:
\[
  XS(k_z, k_x)
  = B \cdot \bigl[B \cdot e^{-j(k_z \Delta z + k_x \Delta x)}\bigr]^{*}
  = |B(k_z, k_x)|^{2} \cdot e^{j(k_z \Delta z + k_x \Delta x)}
\]

Extracting the phase via \texttt{np.angle(XS)}, the real amplitude factor
$|B|^2$ vanishes, leaving:
\[
  \boxed{\Phi(k_z, k_x) = k_z\,\Delta z + k_x\,\Delta x}
\]

\begin{itemize}
  \item This is a \textbf{flat plane} passing through the origin $(0,0,0)$.
  \item Slope along $k_z$ $\;\Rightarrow\;$ vertical displacement $\Delta z$.
  \item Slope along $k_x$ $\;\Rightarrow\;$ horizontal displacement $\Delta x$.
\end{itemize}

\end{frame}

% ─── Frame 3 ─────────────────────────────────────────────────────────────────
\begin{frame}{3. The Overdetermined System — Least Squares Plane Fit}

A discrete GPR image yields \textbf{thousands} of frequency-domain pixels.
Each valid bin $i = 1, \dots, N$ provides an independent equation:
\[
  \Phi_i = k_{z,i}\,\Delta z + k_{x,i}\,\Delta x + c
\]
where $c$ absorbs minor calibration biases.

\bigskip
This is cast as the overdetermined matrix problem $A\,\mathbf{u} = \bm{\Phi}$:
\[
  \begin{bmatrix}
    k_{z,1} & k_{x,1} & 1 \\
    k_{z,2} & k_{x,2} & 1 \\
    \vdots  & \vdots  & \vdots \\
    k_{z,N} & k_{x,N} & 1
  \end{bmatrix}
  \begin{bmatrix} \Delta z \\ \Delta x \\ c \end{bmatrix}
  =
  \begin{bmatrix} \Phi_1 \\ \Phi_2 \\ \vdots \\ \Phi_N \end{bmatrix}
\]

\medskip
\texttt{np.linalg.lstsq} solves this via SVD, compressing thousands of
measurements into two precise scalars: $\Delta z$ and $\Delta x$.

\end{frame}

% ─── Frame 4a ────────────────────────────────────────────────────────────────
\begin{frame}{4. Why Weights and the Mask are Crucial}

Solving the matrix naively causes error spikes: frequency bins outside the
antenna band contain only \emph{random numerical noise}, which destabilises
the fit.

Two mathematical safeguards enforce sub-wavelength accuracy:

\bigskip
\textbf{A. The Band-Pass Mask}

Restricting the matrix to $|k_z|, |k_x| < 1.4\,k_{zc}$ keeps the algorithm
inside the coherent envelope of the Ricker wavelet pulse.

\medskip
This guarantees that the total phase rotation $\Phi_i$ never crosses
$\pm\pi$ radians — preventing phase \emph{wrapping}.

\end{frame}

% ─── Frame 4b ────────────────────────────────────────────────────────────────
\begin{frame}{4. Why Weights and the Mask are Crucial (cont.)}

\textbf{B. Amplitude Weighting ($W$)}

The linear system is pre-multiplied by a diagonal weight matrix $W$, where
each weight $W_i = |XS|_i$ is the spectral magnitude:
\[
  W A\,\mathbf{u} = W\,\bm{\Phi}
\]

The optimisation goal becomes minimising the \textbf{energy-weighted} residual:
\[
  \min \sum_{i=1}^{N}
    |XS_i|^{2} \cdot
    \Bigl(\Phi_i - \bigl(k_{z,i}\,\Delta z + k_{x,i}\,\Delta x + c\bigr)\Bigr)^{2}
\]

\begin{itemize}
  \item Frequency bins at \textbf{peak antenna power} carry maximum weight.
  \item Bins near the \textbf{noise floor} are mathematically suppressed.
\end{itemize}

\medskip
The result: two extraordinarily precise displacement estimates,
$\Delta z$ and $\Delta x$.

\end{frame}

% ─── Application: phase-shift method results ─────────────────────────────────
\begin{frame}{Results: Horizontal, Vertical \& Diagonal Shift — $\tfrac{1}{4}\lambda$}
\begin{center}
  \includegraphics[width=\linewidth,height=0.80\textheight,keepaspectratio]{Phase_shift_method_test_0p25_lambda.png}
\end{center}
\end{frame}

\begin{frame}{Results: Horizontal, Vertical \& Diagonal Shift — $\tfrac{1}{8}\lambda$}
\begin{center}
  \includegraphics[width=\linewidth,height=0.80\textheight,keepaspectratio]{Phase_shift_method_test_0p125_lambda.png}
\end{center}
\end{frame}

\begin{frame}{Results: Horizontal, Vertical \& Diagonal Shift — $\tfrac{1}{16}\lambda$}
\begin{center}
  \includegraphics[width=\linewidth,height=0.80\textheight,keepaspectratio]{Phase_shift_method_test_0p0625_lambda.png}
\end{center}
\end{frame}

\begin{frame}{Results: Horizontal, Vertical \& Diagonal Shift — $\tfrac{1}{32}\lambda$}
\begin{center}
  \includegraphics[width=\linewidth,height=0.80\textheight,keepaspectratio]{Phase_shift_method_test_0p03125_lambda.png}
\end{center}
\end{frame}

% ─── Application: gprMax data ────────────────────────────────────────────────
\begin{frame}{gprMax Data — $\tfrac{1}{4}\lambda$}
\begin{center}
  \includegraphics[width=\linewidth,height=0.80\textheight,keepaspectratio]{gprMax_test_0p25_lambda.png}
\end{center}
\end{frame}

\begin{frame}{gprMax Data — $\tfrac{1}{8}\lambda$}
\begin{center}
  \includegraphics[width=\linewidth,height=0.80\textheight,keepaspectratio]{gprMax_test_0p125_lambda.png}
\end{center}
\end{frame}

\begin{frame}{gprMax Data — $\tfrac{1}{16}\lambda$}
\begin{center}
  \includegraphics[width=\linewidth,height=0.80\textheight,keepaspectratio]{gprMax_test_0p0625_lambda.png}
\end{center}
\end{frame}

\end{document}
